% chapters/bab1-pendahuluan.tex -- BAB 1 PENDAHULUAN
% Contoh teks ilmiah di bawah ini adalah placeholder untuk memperlihatkan
% hasil format (thesis.md); ganti dengan isi penelitian yang sebenarnya.

\chapter{PENDAHULUAN}

\section{Latar Belakang Masalah}
Perkembangan teknologi informasi telah mendorong lonjakan volume data
tekstual yang dihasilkan pengguna internet, terutama dalam bentuk
ulasan produk pada platform perdagangan elektronik. Ulasan tersebut
memuat opini konsumen yang sangat berharga bagi pelaku usaha, baik
untuk mengevaluasi kualitas produk maupun untuk merumuskan strategi
pemasaran. Namun demikian, jumlah ulasan yang terus bertambah setiap
hari membuat proses pembacaan dan penilaian opini secara manual
menjadi tidak efisien, sehingga dibutuhkan suatu pendekatan otomatis
yang mampu mengklasifikasikan sentimen ulasan tersebut ke dalam
kategori positif, negatif, atau netral.

Berbagai pendekatan pembelajaran mesin telah digunakan untuk
menyelesaikan permasalahan klasifikasi sentimen, mulai dari metode
statistik konvensional seperti Naive Bayes dan Support Vector Machine,
hingga metode pembelajaran mendalam (deep learning) yang mampu
menangkap hubungan kontekstual antar kata dalam suatu kalimat. Salah
satu arsitektur yang banyak digunakan untuk data sekuensial seperti
teks adalah Long Short-Term Memory (LSTM), sebuah varian dari
Recurrent Neural Network (RNN) yang dirancang untuk mengatasi masalah
vanishing gradient pada data berurutan yang panjang.

Bahasa Indonesia memiliki karakteristik morfologis dan tata bahasa
yang berbeda dari bahasa Inggris, seperti penggunaan imbuhan, kata
tidak baku, dan campur kode dengan bahasa daerah maupun bahasa asing
pada ulasan produk. Karakteristik ini menjadi tantangan tersendiri
dalam praproses data sebelum data tersebut dapat digunakan untuk
melatih model klasifikasi. Oleh karena itu, penelitian ini berfokus
pada penerapan metode LSTM untuk klasifikasi sentimen ulasan produk
berbahasa Indonesia, dengan memperhatikan tahapan praproses yang
sesuai dengan karakteristik bahasa tersebut.

\section{Rumusan Masalah}
Berdasarkan uraian latar belakang di atas, rumusan masalah dalam
penelitian ini adalah sebagai berikut.
\begin{enumerate}
  \item Bagaimana tahapan praproses data yang tepat untuk ulasan
        produk berbahasa Indonesia sebelum digunakan pada model LSTM.
  \item Bagaimana arsitektur model LSTM yang optimal untuk
        mengklasifikasikan sentimen ulasan produk berbahasa Indonesia.
  \item Bagaimana tingkat akurasi model LSTM dibandingkan dengan
        metode pembelajaran mesin konvensional pada dataset yang sama.
\end{enumerate}

\subsection{Batasan masalah penelitian}
Agar penelitian ini lebih terarah, permasalahan dibatasi pada
beberapa hal berikut: data yang digunakan adalah ulasan produk
berbahasa Indonesia yang diambil dari salah satu platform perdagangan
elektronik; klasifikasi sentimen dibatasi pada tiga kelas, yaitu
positif, negatif, dan netral; serta model yang dibangun tidak
dibandingkan dengan model berbasis Transformer mengingat keterbatasan
sumber daya komputasi yang tersedia.

\section{Tujuan Penelitian}
Penelitian ini bertujuan untuk merancang, membangun, dan mengevaluasi
model klasifikasi sentimen berbasis LSTM untuk ulasan produk
berbahasa Indonesia, serta membandingkan kinerjanya dengan metode
pembelajaran mesin konvensional sebagai baseline.

\section{Manfaat Penelitian}
Manfaat penelitian ini dijabarkan ke dalam dua aspek, yaitu manfaat
teoretis dan manfaat praktis.

\subsection{Manfaat teoretis bagi pengembangan ilmu}
Secara teoretis, penelitian ini diharapkan dapat memperkaya kajian di
bidang pemrosesan bahasa alami, khususnya mengenai penerapan metode
LSTM pada data berbahasa Indonesia yang memiliki karakteristik
morfologis khas.

\subsection{Manfaat praktis bagi pelaku usaha}
Secara praktis, hasil penelitian ini diharapkan dapat dimanfaatkan
oleh pelaku usaha untuk memantau opini konsumen secara otomatis dan
efisien, sehingga dapat digunakan sebagai bahan pertimbangan dalam
pengambilan keputusan bisnis.
